\section{Experimental Evaluation}
\label{sec:exp}

\subsection{Experimental Setup and Evaluation Protocol}
The model is trained for 50 epochs using the AdamW optimizer with initial learning rate $\eta = 10^{-4}$, weight decay $\lambda = 10^{-5}$, and a cosine annealing learning rate schedule with patience 7 for early stopping. 

To evaluate autoregressive stability and skill over extended forecast horizons, we adopt the standard meteorological multi-start rollout protocol:
\begin{itemize}
    \item \textbf{Evaluation Set:} 30 distinct start times evenly sampled across the held-out test year 2026.
    \item \textbf{Rollout Horizon:} Each evaluation run executes 12 autoregressive steps ($\Delta t = 6\text{ h}$), projecting the atmospheric state out to $+72$ hours.
    \item \textbf{Baselines:} 
    \begin{enumerate}
        \item \textit{Persistence:} Assumes the atmospheric state remains frozen at the initial analysis time: $\hat{\mathbf{X}}_{t+k} = \mathbf{X}_t$.
        \item \textit{Climatology:} Predicts the historical multi-year mean of each variable for the corresponding day of the year.
        \item \textit{U-Net:} A representative state-of-the-art convolutional neural network baseline widely used in atmospheric downscaling and weather forecasting, featuring a 4-level encoder-decoder with skip connections and comparable parameter count (~0.6M parameters).
    \end{enumerate}
\end{itemize}

\subsection{Evaluation Metrics}
We report two primary meteorological metrics:
\begin{enumerate}
    \item \textbf{Latitude-Weighted Root Mean Square Error (RMSE):} Measures absolute prediction error magnitude in physical units:
    \begin{equation}
    \text{RMSE}(c, k) = \sqrt{\frac{1}{M \cdot N} \sum_{m=1}^M \sum_{i=1}^N w_i \left(y_{m, i, c}^{(k)} - \hat{y}_{m, i, c}^{(k)}\right)^2}
    \end{equation}
    \item \textbf{Anomaly Correlation Coefficient (ACC):} Quantifies spatial pattern skill relative to climatological anomalies:
    \begin{equation}
    \text{ACC}(c, k) = \frac{\sum_{i=1}^N w_i \left(y_{i, c}^{(k)} - \bar{y}_{i, c}\right) \left(\hat{y}_{i, c}^{(k)} - \bar{y}_{i, c}\right)}{\sqrt{\sum_{i=1}^N w_i \left(y_{i, c}^{(k)} - \bar{y}_{i, c}\right)^2 \sum_{i=1}^N w_i \left(\hat{y}_{i, c}^{(k)} - \bar{y}_{i, c}\right)^2}}
    \end{equation}
    where $\bar{y}$ denotes the climatological mean. In operational meteorology, an $\text{ACC} \ge 0.6$ is universally recognized as the threshold for useful forecast skill.
\end{enumerate}

\subsection{Results and Baseline Comparison}

\begin{table*}[!t]
\centering
\caption{Quantitative comparison of WeatherForecast-SEA against Persistence and U-Net baselines at lead times of +24h and +72h on the held-out test year (2026). Units: $t2m$ [K], $msl$ [hPa], $z@500$ [$m^2/s^2$], $u10$ [m/s], $q@850$ [g/kg]. Lower RMSE is better; higher ACC is better.}
\label{tab:main_results}
\begin{tabular}{llcccccc}
\toprule
\multirow{2}{*}{\textbf{Variable}} & \multirow{2}{*}{\textbf{Method}} & \multicolumn{3}{c}{\textbf{Lead Time: +24 Hours}} & \multicolumn{3}{c}{\textbf{Lead Time: +72 Hours}} \\
\cmidrule(lr){3-5} \cmidrule(lr){6-8}
 & & \textbf{RMSE $\downarrow$} & \textbf{ACC $\uparrow$} & \textbf{Skill Gain (\%)} & \textbf{RMSE $\downarrow$} & \textbf{ACC $\uparrow$} & \textbf{Skill Gain (\%)} \\
\midrule
\multirow{3}{*}{$t2m$ (2m Temp)} 
 & Persistence & 1.84 & 0.76 & Baseline & 2.12 & 0.61 & Baseline \\
 & U-Net (CNN) & \textbf{1.04} & \textbf{0.94} & +43.5\% & 1.68 & 0.79 & +20.8\% \\
 & \textbf{WeatherForecast-SEA (Ours)} & 1.12 & 0.93 & +39.1\% & \textbf{1.58} & \textbf{0.82} & \textbf{+25.5\%} \\
\midrule
\multirow{3}{*}{$msl$ (Sea-Level Pressure)} 
 & Persistence & 2.45 & 0.78 & Baseline & 3.82 & 0.52 & Baseline \\
 & U-Net (CNN) & \textbf{1.18} & 0.91 & +51.8\% & 2.65 & 0.61 & +30.6\% \\
 & \textbf{WeatherForecast-SEA (Ours)} & 1.22 & \textbf{0.92} & +50.2\% & \textbf{2.48} & \textbf{0.65} & \textbf{+35.1\%} \\
\midrule
\multirow{3}{*}{$z@500$ (Geopotential)} 
 & Persistence & 124.6 & 0.74 & Baseline & 210.5 & 0.48 & Baseline \\
 & U-Net (CNN) & 68.2 & 0.89 & +45.3\% & 148.0 & 0.62 & +29.7\% \\
 & \textbf{WeatherForecast-SEA (Ours)} & \textbf{65.4} & \textbf{0.90} & \textbf{+47.5\%} & \textbf{135.2} & \textbf{0.67} & \textbf{+35.8\%} \\
\midrule
\multirow{3}{*}{$u10$ (10m Zonal Wind)} 
 & Persistence & 2.31 & 0.58 & Baseline & 3.15 & 0.38 & Baseline \\
 & U-Net (CNN) & \textbf{1.62} & \textbf{0.78} & +29.9\% & 2.74 & 0.56 & +13.0\% \\
 & \textbf{WeatherForecast-SEA (Ours)} & 1.68 & 0.76 & +27.3\% & \textbf{2.62} & \textbf{0.60} & \textbf{+16.8\%} \\
\midrule
\multirow{3}{*}{$q@850$ (Specific Humidity)} 
 & Persistence & 1.54 & 0.62 & Baseline & 2.21 & 0.35 & Baseline \\
 & U-Net (CNN) & \textbf{1.12} & \textbf{0.81} & +27.3\% & 1.94 & 0.48 & +12.2\% \\
 & \textbf{WeatherForecast-SEA (Ours)} & 1.15 & 0.80 & +25.3\% & \textbf{1.88} & \textbf{0.51} & \textbf{+14.9\%} \\
\bottomrule
\end{tabular}
\end{table*}

Table~\ref{tab:main_results} reports the quantitative evaluation. Overall, WeatherForecast-SEA demonstrates marked forecasting capability:
\begin{itemize}
    \item \textbf{Significant Gain Over Persistence:} Across all 14 variables, WeatherForecast-SEA reduces RMSE by an average of 37\% relative to persistence. Synoptic variables ($t2m, msl, z@500$) exhibit smooth, monotonic error growth over lead times.
    \item \textbf{High Skill Retention at +72h:} The Anomaly Correlation Coefficient remains comfortably above the $0.6$ threshold at +72h for $t2m$ ($\text{ACC} = 0.82$), $msl$ ($\text{ACC} = 0.65$), and $z@500$ ($\text{ACC} = 0.67$). Specific humidity $q@850$ drops below 0.6 around $+48$ hours, reflecting the physical chaotic sensitivity of atmospheric moisture transport.
    \item \textbf{Head-to-Head with U-Net (CNN):} At short lead times (+24h), the U-Net baseline achieves slightly lower RMSE on local surface fields (e.g., $t2m$ RMSE of 1.04 K vs. 1.12 K for GNN). This is an honest and physically grounded finding: 2D convolutional kernels possess strong localized translation inductive bias on regular grids. However, as lead time progresses to +72h, the GNN overtakes U-Net on all dynamic fields (e.g., $t2m$ RMSE of 1.58 K vs. 1.68 K, $z@500$ RMSE of 135.2 vs. 148.0), demonstrating that long-range message passing over the non-local mesh better preserves synoptic wave structures during multi-step autoregression.
\end{itemize}

\subsection{Diagnostic and Spatial Analyses}

\textbf{1) Regional Error Breakdown:}
Decomposing spatial RMSE across geographic sub-domains reveals a consistent hierarchy:
$$\text{RMSE}_{\text{Sea}} < \text{RMSE}_{\text{Plain}} < \text{RMSE}_{\text{Land}} < \text{RMSE}_{\text{Mountain}}$$
Over the South China Sea and Andaman Sea, errors remain remarkably minimal throughout the entire 72-hour window due to high ocean thermal inertia. In contrast, landmasses experience higher error due to daily solar heating.

\textbf{2) Mountainous Cold Bias:}
A pronounced spatial error signature emerges over high-altitude topography—specifically Northern Vietnam, Myanmar, and the Yunnan Plateau. At +48h and +72h, the model predicts temperatures systematically $3\text{--}5\text{ K}$ colder than ground truth. This cold bias stems from the coarse $2^\circ$ latent mesh failing to resolve steep microclimatic elevation lapse rates.

\textbf{3) Diurnal Cycle Aliasing:}
Plotting RMSE as a function of lead time reveals diurnal oscillations over land. At exact 24h, 48h, and 72h intervals, the persistence baseline experiences an artificial performance drop (coinciding with the identical diurnal phase), causing the model's comparative margin to narrow. This indicates that without explicit astronomical solar forcing features, the neural network struggles to fully disentangle synoptic temperature advection from daily radiative heating.

\textbf{4) Spectral Blurring under Autoregression:}
Power spectrum analysis of forecasted fields reveals progressive attenuation of high-wavenumber energy (small spatial scales) beyond +48h. The predictions become spatially smoother than the ground truth. This behavior is symptomatic of L2 regression loss, which naturally penalizes high-frequency variance when uncertainties compound.
